\documentclass[12pt,a4paper]{article}
 
% ─── Package ───────────────────────────────────────────────────────────────────
\usepackage[utf8]{inputenc}
\usepackage[T1]{fontenc}
\usepackage{amsmath, amssymb}
\usepackage{geometry}
\usepackage{graphicx}
\usepackage{xcolor}
\usepackage{listings}
\usepackage{hyperref}
\usepackage{booktabs}
\usepackage{enumitem}
\usepackage{tcolorbox}
\usepackage{mdframed}
\usepackage{fancyhdr}
\usepackage{titlesec}
 
\geometry{margin=2.5cm}
 
% ─── Warna Tema ────────────────────────────────────────────────────────────────
\definecolor{codebg}{RGB}{245,247,250}
\definecolor{codeframe}{RGB}{180,195,215}
\definecolor{keywordcolor}{RGB}{0,112,192}
\definecolor{commentcolor}{RGB}{100,160,100}
\definecolor{stringcolor}{RGB}{180,70,50}
\definecolor{notebg}{RGB}{255,249,230}
\definecolor{noteframe}{RGB}{230,180,50}
\definecolor{tipbg}{RGB}{230,245,235}
\definecolor{tipframe}{RGB}{50,160,90}
\definecolor{sectioncolor}{RGB}{30,80,150}
 
% ─── Style Section ─────────────────────────────────────────────────────────────
\titleformat{\section}{\Large\bfseries\color{sectioncolor}}{}{0em}{\thesection\quad}[\titlerule]
\titleformat{\subsection}{\large\bfseries\color{sectioncolor!80}}{}{0em}{\thesubsection\quad}
 
% ─── Listings (Python) ─────────────────────────────────────────────────────────
\lstdefinestyle{pythonstyle}{
  language=Python,
  backgroundcolor=\color{codebg},
  frame=single,
  rulecolor=\color{codeframe},
  framerule=0.8pt,
  basicstyle=\ttfamily\small,
  keywordstyle=\bfseries\color{keywordcolor},
  commentstyle=\itshape\color{commentcolor},
  stringstyle=\color{stringcolor},
  showstringspaces=false,
  breaklines=true,
  numbers=left,
  numberstyle=\tiny\color{gray},
  numbersep=8pt,
  tabsize=4,
  captionpos=b,
}
\lstset{style=pythonstyle}
 
% ─── Kotak Catatan & Tips ──────────────────────────────────────────────────────
\newmdenv[
  backgroundcolor=notebg,
  linecolor=noteframe,
  linewidth=2pt,
  leftline=true, topline=false, rightline=false, bottomline=false,
  innerleftmargin=10pt, innerrightmargin=10pt,
  innertopmargin=6pt, innerbottommargin=6pt,
  skipabove=8pt, skipbelow=8pt
]{notebox}
 
\newmdenv[
  backgroundcolor=tipbg,
  linecolor=tipframe,
  linewidth=2pt,
  leftline=true, topline=false, rightline=false, bottomline=false,
  innerleftmargin=10pt, innerrightmargin=10pt,
  innertopmargin=6pt, innerbottommargin=6pt,
  skipabove=8pt, skipbelow=8pt
]{tipbox}
 
% ─── Header / Footer ───────────────────────────────────────────────────────────
\pagestyle{fancy}
\fancyhf{}
\fancyhead[L]{\small\textit{Pengolahan Citra Digital -- Bab 2}}
\fancyhead[R]{\small\textit{Dasar-Dasar Pembentukan Citra}}
\fancyfoot[C]{\thepage}
\renewcommand{\headrulewidth}{0.4pt}
 
% ─── Judul ─────────────────────────────────────────────────────────────────────
\title{
  {\LARGE\bfseries Pengolahan Citra Digital}\\[0.5em]
  {\large Bab 2: Dasar-Dasar Pembentukan Citra}\\[0.3em]
  {\normalsize Catatan Kuliah dengan Penjelasan \& Contoh Kode Python}
}
\author{}
\date{}
 
% ===============================================================================
\begin{document}
\maketitle
\thispagestyle{fancy}
\tableofcontents
\newpage
 
% ===============================================================================
\section{Persepsi Visual Manusia}
\label{sec:visual}
 
\subsection{Kenapa Perlu Belajar Tentang Mata?}
 
Mungkin kamu bertanya-tanya, \textit{``Ini kan kuliah teknik, ngapain bahas mata manusia?''}
Alasannya cukup masuk akal: sebagian besar citra digital yang kita proses pada akhirnya
dikonsumsi oleh manusia. Jadi kalau kita tidak paham cara kerja mata, bisa jadi algoritma
yang kita buat terlalu berlebihan atau malah kurang pas.
 
Bayangkan kamu sibuk \textit{sharpening} (mempertajam) gambar selama berjam-jam, padahal
mata manusia sudah punya mekanisme penajaman alaminya sendiri.
 
\subsection{Adaptasi dan Diskriminasi Kecerahan}
 
Mata manusia itu fleksibel banget. Berbeda dengan kamera yang punya rentang dinamis
(\textit{dynamic range}) tetap, mata kita terus-menerus menyesuaikan diri dengan tingkat
cahaya rata-rata di sekitar kita, yang biasa disebut \textbf{level adaptasi} $B_a$.
 
\begin{itemize}[leftmargin=*]
  \item \textbf{Rentang total:} Sangat luas -- dari cahaya bintang yang redup sampai
        sinar matahari terik.
  \item \textbf{Rentang pada satu waktu:} Hanya sebagian kecil di sekitar $B_a$. Makanya
        kita susah melihat detail di area yang sangat gelap \textit{dan} sangat terang
        sekaligus.
\end{itemize}
 
\subsection{Fenomena Visual Penting}
 
\subsubsection{Mach Bands}
 
Kalau ada dua area dengan intensitas berbeda yang bersebelahan, mata kita akan
``merasakan'' garis terang atau gelap \textit{semu} di perbatasannya -- padahal secara
fisik intensitasnya konstan. Fenomena ini disebut \textbf{Mach Bands}.
 
\begin{tipbox}
  \textbf{Implikasi praktis:} Karena mata sudah ``mempertajam'' perbatasan secara alami,
  kita perlu hati-hati saat menerapkan algoritma \textit{sharpening} berlebihan. Hasilnya
  bisa terlihat tidak natural.
\end{tipbox}
 
\subsubsection{Simultaneous Contrast}
 
Persepsi kecerahan suatu objek sangat dipengaruhi oleh latar belakangnya. Kotak
abu-abu yang sama persis akan terlihat \textit{lebih gelap} jika ditempatkan di atas
latar putih, dibandingkan jika ditempatkan di atas latar hitam.
 
\begin{notebox}
  \textbf{Pertanyaan Reflektif:} Jika mata sudah melakukan ``penajaman alami'' seperti
  Mach Bands, apakah kita perlu mempertajam citra digital secara agresif lewat algoritma?
  Tidak selalu -- tergantung konteks dan tujuan penggunaannya.
\end{notebox}
 
% ===============================================================================
\section{Cahaya dan Spektrum Elektromagnetik}
\label{sec:em}
 
\subsection{Citra sebagai Energi Elektromagnetik}
 
Citra yang kita olah umumnya berasal dari energi elektromagnetik (EM) -- bisa cahaya
tampak, inframerah, sinar-X, dan lain-lain. Pemahaman tentang sifat EM penting agar
kita tahu \textit{apa yang sebenarnya kita ukur}.
 
\subsection{Sifat Gelombang}
 
Hubungan fundamental antara panjang gelombang $\lambda$, frekuensi $\nu$, dan kecepatan
cahaya $c$:
 
\begin{equation}
  \lambda \nu = c
  \label{eq:wave}
\end{equation}
 
di mana $c \approx 3 \times 10^8\ \text{m/s}$. Cahaya tampak oleh manusia berada pada
rentang $\lambda \approx 400\text{--}700\ \text{nm}$.
 
\subsection{Sifat Partikel (Foton)}
 
Dalam model kuantum, energi satu foton $E$ berbanding lurus dengan frekuensinya:
 
\begin{equation}
  E = h\nu
  \label{eq:photon}
\end{equation}
 
di mana $h = 6.626 \times 10^{-34}\ \text{J}\cdot\text{s}$ adalah konstanta Planck.
 
\begin{notebox}
  \textbf{Intuisi:} Sinar-X punya frekuensi jauh lebih tinggi dari cahaya tampak, sehingga
  energi fotonnya lebih besar. Inilah kenapa sinar-X bisa menembus jaringan tubuh, sementara
  cahaya tampak tidak bisa.
\end{notebox}
 
% ===============================================================================
\section{Akuisisi Citra dan Model Pembentukan Citra}
\label{sec:akuisisi}
 
\subsection{Bagaimana Citra Terbentuk?}
 
Secara sederhana, proses pembentukan citra adalah:
\begin{enumerate}
  \item Sumber cahaya memancarkan energi ke pemandangan (\textit{scene}).
  \item Objek memantulkan atau meneruskan cahaya tersebut.
  \item Sensor (misalnya CCD pada kamera digital) menangkap energi yang sampai.
  \item Hasilnya direkam sebagai citra digital.
\end{enumerate}
 
\subsection{Model Matematis: Fungsi Citra}
 
Sebuah citra didefinisikan sebagai fungsi dua dimensi:
 
\begin{equation}
  f(x, y) = i(x, y) \cdot r(x, y)
  \label{eq:image_model}
\end{equation}
 
\noindent di mana:
\begin{itemize}
  \item $i(x,y)$ -- \textbf{Iluminasi} (\textit{illumination}): jumlah cahaya yang jatuh
        pada titik $(x,y)$. Nilainya selalu positif:
        \[0 < i(x,y) < \infty\]
  \item $r(x,y)$ -- \textbf{Reflektansi} (\textit{reflectance}): proporsi cahaya yang
        dipantulkan oleh objek. Nilainya antara 0 dan 1:
        \[0 \leq r(x,y) \leq 1\]
        Nilai $r = 0$ berarti objek menyerap semua cahaya (hitam sempurna), sedangkan
        $r = 1$ berarti memantulkan semua cahaya (putih sempurna).
\end{itemize}
 
\begin{tipbox}
  \textbf{Analogi gampang:} Bayangkan kamu menyorotkan senter ke dinding.
  $i(x,y)$ adalah seberapa kuat senternya, sedangkan $r(x,y)$ adalah seberapa
  ``putih'' dindingnya. Hasil yang kamera tangkap adalah perkalian keduanya -- $f(x,y)$.
\end{tipbox}
 
% ===============================================================================
\section{Sampling dan Kuantisasi}
\label{sec:sampling}
 
\subsection{Dari Dunia Nyata ke Komputer}
 
Dunia nyata itu kontinu -- intensitas cahaya bisa bernilai apa saja di sembarang titik.
Komputer, di sisi lain, hanya bisa menangani bilangan diskret. Di sinilah peran
\textit{sampling} dan \textit{kuantisasi}.
 
\begin{itemize}
  \item \textbf{Sampling} -- Digitalisasi koordinat spasial $(x, y)$. Menentukan berapa
        banyak pixel yang ada (resolusi spasial).
  \item \textbf{Kuantisasi} -- Digitalisasi nilai intensitas $f$. Menentukan berapa banyak
        tingkat keabuan yang tersedia (resolusi intensitas).
\end{itemize}
 
\subsection{Representasi Matriks}
 
Sebuah citra berukuran $M \times N$ direpresentasikan sebagai matriks:
 
\begin{equation}
  \mathbf{A} =
  \begin{pmatrix}
    a_{0,0}   & a_{0,1}   & \cdots & a_{0,N-1}   \\
    a_{1,0}   & a_{1,1}   & \cdots & a_{1,N-1}   \\
    \vdots    & \vdots    & \ddots & \vdots      \\
    a_{M-1,0} & a_{M-1,1} & \cdots & a_{M-1,N-1}
  \end{pmatrix}
  \label{eq:matrix}
\end{equation}
 
Setiap elemen $a_{i,j}$ disebut \textbf{pixel} (\textit{picture element}).
 
\subsection{Kebutuhan Penyimpanan}
 
Jika citra memiliki $L = 2^k$ tingkat intensitas, maka jumlah bit yang
dibutuhkan untuk menyimpan citra $M \times N$ adalah:
 
\begin{equation}
  b = M \times N \times k
  \label{eq:storage}
\end{equation}
 
\begin{notebox}
  \textbf{Contoh:} Citra \textit{grayscale} 8-bit ($k=8$, artinya $L=256$ level) berukuran
  $1920 \times 1080$ membutuhkan:
  \[b = 1920 \times 1080 \times 8 = 16{,}588{,}800\ \text{bit} \approx 1.98\ \text{MB}\]
  Ini adalah ukuran sebelum kompresi. Format JPEG atau PNG mengompresi lebih lanjut.
\end{notebox}
 
\subsection{Contoh Kode Python}
 
\begin{lstlisting}[caption={Membaca citra dan memeriksa properti dasarnya},
                   label={lst:basic}]
import numpy as np
from PIL import Image
import matplotlib.pyplot as plt
 
# Baca citra dan konversi ke grayscale
img = Image.open("contoh.jpg").convert("L")   # "L" = grayscale
A   = np.array(img)                           # Konversi ke NumPy array
 
M, N = A.shape
k    = 8                                      # Bit-depth (asumsi 8-bit)
L    = 2**k                                   # Jumlah level: 256
 
print(f"Ukuran citra  : {M} x {N} pixel")
print(f"Bit-depth     : {k} bit  =>  L = {L} tingkat keabuan")
print(f"Kebutuhan bit : {M*N*k:,} bit  ({M*N*k/8/1024:.1f} KB)")
print(f"Nilai min     : {A.min()}")
print(f"Nilai max     : {A.max()}")
 
# Tampilkan citra dan histogram intensitas
fig, axes = plt.subplots(1, 2, figsize=(10, 4))
 
axes[0].imshow(A, cmap="gray", vmin=0, vmax=255)
axes[0].set_title("Citra Grayscale")
axes[0].axis("off")
 
axes[1].hist(A.ravel(), bins=256, range=(0, 255),
             color="steelblue", edgecolor="none")
axes[1].set_title("Histogram Intensitas")
axes[1].set_xlabel("Nilai Pixel")
axes[1].set_ylabel("Frekuensi")
 
plt.tight_layout()
plt.show()
\end{lstlisting}
 
% ===============================================================================
\section{Hubungan Antar Pixel}
\label{sec:hubungan}
 
\subsection{Konsep Ketetanggaan (\textit{Neighborhood})}
 
Untuk berbagai operasi seperti segmentasi atau deteksi tepi, kita perlu tahu pixel mana
yang dianggap ``bertetangga'' dengan pixel $p$ di koordinat $(x, y)$.
 
\subsubsection{4-Tetangga: $N_4(p)$}
 
Hanya pixel di atas, bawah, kiri, dan kanan:
\begin{equation}
  N_4(p) = \{(x+1,y),\ (x-1,y),\ (x,y+1),\ (x,y-1)\}
\end{equation}
 
\subsubsection{8-Tetangga: $N_8(p)$}
 
$N_4(p)$ ditambah keempat tetangga diagonal $N_D(p)$:
\begin{equation}
  N_8(p) = N_4(p) \cup N_D(p)
\end{equation}
 
\subsection{Jenis Ketertetanggaan (\textit{Adjacency})}
 
\begin{itemize}
  \item \textbf{4-adjacency:} Dua pixel bertetangga jika salah satunya ada di $N_4$
        pixel lainnya.
  \item \textbf{8-adjacency:} Dua pixel bertetangga jika salah satunya ada di $N_8$
        pixel lainnya.
  \item \textbf{m-adjacency} (mixed): Kombinasi 4- dan 8-adjacency untuk menghilangkan
        ambiguitas jalur. Berguna saat definisi ``path'' perlu unik.
\end{itemize}
 
\subsection{Ukuran Jarak}
 
Diberikan dua pixel $p = (x, y)$ dan $q = (u, v)$:
 
\begin{align}
  D_e(p,q) &= \sqrt{(x-u)^2 + (y-v)^2}
  \quad \text{(Euclidean)} \label{eq:eucl} \\[6pt]
  D_4(p,q) &= |x-u| + |y-v|
  \quad \text{(City-block / Manhattan)} \label{eq:city} \\[6pt]
  D_8(p,q) &= \max\bigl(|x-u|,\ |y-v|\bigr)
  \quad \text{(Chessboard)} \label{eq:chess}
\end{align}
 
\begin{notebox}
  \textbf{Analogi:}
  \begin{itemize}[nosep]
    \item $D_e$ -- jarak garis lurus seperti burung terbang.
    \item $D_4$ -- jarak berjalan di kota Manhattan yang jalannya berbentuk grid.
    \item $D_8$ -- jarak yang ditempuh raja di papan catur.
  \end{itemize}
\end{notebox}
 
\subsection{Contoh Kode Python}
 
\begin{lstlisting}[caption={Menghitung tiga jenis jarak antar pixel},
                   label={lst:distance}]
import numpy as np
import matplotlib.pyplot as plt
 
def euclidean(p, q):
    """Jarak Euclidean -- garis lurus."""
    return np.sqrt((p[0]-q[0])**2 + (p[1]-q[1])**2)
 
def city_block(p, q):
    """City-block (D4) -- berjalan di grid Manhattan."""
    return abs(p[0]-q[0]) + abs(p[1]-q[1])
 
def chessboard(p, q):
    """Chessboard (D8) -- gerakan raja di catur."""
    return max(abs(p[0]-q[0]), abs(p[1]-q[1]))
 
# Uji coba
p = (0, 0)
q = (3, 4)
print(f"Pixel p = {p},  pixel q = {q}")
print(f"D_euclidean = {euclidean(p, q):.4f}")   # => 5.0000
print(f"D_4         = {city_block(p, q)}")       # => 7
print(f"D_8         = {chessboard(p, q)}")        # => 4
 
# Visualisasi area dengan jarak <= 2 dari pusat
center = (5, 5)
radius = 2
size   = 11
 
fig, axes = plt.subplots(1, 2, figsize=(10, 5))
titles = ["D4 (City-block) <= 2", "D8 (Chessboard) <= 2"]
 
for ax, dist_fn, title in zip(axes, [city_block, chessboard], titles):
    grid = np.zeros((size, size))
    for r in range(size):
        for c in range(size):
            if dist_fn(center, (r, c)) <= radius:
                grid[r, c] = 1
    ax.imshow(grid, cmap="Blues", vmin=0, vmax=1.5)
    ax.plot(center[1], center[0], "ro", markersize=8, label="Pusat")
    ax.set_title(title)
    ax.legend()
 
plt.tight_layout()
plt.show()
\end{lstlisting}
 
% ===============================================================================
\section{Perangkat Matematika Utama}
\label{sec:math}
 
\subsection{Operasi Elementwise}
 
Hampir semua operasi pada citra bersifat \textit{elementwise} -- dilakukan
pixel demi pixel secara independen. Misalnya, menjumlahkan dua citra:
 
\begin{equation}
  g(x,y) = f_1(x,y) + f_2(x,y) \quad \forall\,(x,y)
\end{equation}
 
Di Python dengan NumPy, operasi ini dilakukan secara otomatis dan efisien lewat mekanisme
\textit{broadcasting}.
 
\subsection{Linearitas dan Prinsip Superposisi}
 
Sebuah operator $H$ disebut \textbf{linear} jika memenuhi prinsip superposisi:
 
\begin{equation}
  H\bigl[a\,f_1 + b\,f_2\bigr] = a\,H[f_1] + b\,H[f_2]
  \label{eq:linear}
\end{equation}
 
Ini penting karena banyak teori pengolahan citra -- termasuk Transformasi Fourier --
mengasumsikan sistem yang linear. Kalau operator tidak linear, banyak teori tersebut
tidak bisa dipakai langsung.
 
\begin{tipbox}
  \textbf{Cek linearitas dengan cepat:} Uji apakah $H[\mathbf{0}] = \mathbf{0}$.
  Jika tidak, operator tersebut pasti tidak linear. Contoh: $H[f] = f + 5$
  tidak linear karena $H[0] = 5 \neq 0$.
\end{tipbox}
 
\subsection{Transformasi Affine (Geometris)}
 
Untuk memanipulasi posisi pixel (rotasi, scaling, translasi), digunakan
\textbf{koordinat homogen}. Transformasi dinyatakan sebagai perkalian matriks $3\times3$:
 
\begin{equation}
  \begin{pmatrix} x' \\ y' \\ 1 \end{pmatrix}
  = \mathbf{A}
  \begin{pmatrix} x \\ y \\ 1 \end{pmatrix}
  \label{eq:affine}
\end{equation}
 
\noindent Beberapa contoh matriks transformasi $\mathbf{A}$:
 
\begin{align}
  \text{Translasi:}\quad
  \mathbf{A}_T &=
  \begin{pmatrix} 1 & 0 & t_x \\ 0 & 1 & t_y \\ 0 & 0 & 1 \end{pmatrix}
  \label{eq:translation} \\[8pt]
  \text{Rotasi sebesar } \theta:\quad
  \mathbf{A}_R &=
  \begin{pmatrix} \cos\theta & -\sin\theta & 0 \\ \sin\theta & \cos\theta & 0 \\ 0 & 0 & 1 \end{pmatrix}
  \label{eq:rotation} \\[8pt]
  \text{Scaling:}\quad
  \mathbf{A}_S &=
  \begin{pmatrix} s_x & 0 & 0 \\ 0 & s_y & 0 \\ 0 & 0 & 1 \end{pmatrix}
  \label{eq:scaling}
\end{align}
 
\noindent Keuntungan koordinat homogen: komposisi transformasi cukup dengan
mengalikan matriks-matriks tersebut:
\[
  \mathbf{A}_{\text{total}} = \mathbf{A}_T \cdot \mathbf{A}_R \cdot \mathbf{A}_S
\]
 
\subsection{Contoh Kode Python}
 
\begin{lstlisting}[caption={Transformasi Affine: rotasi, translasi, dan scaling},
                   label={lst:affine}]
import numpy as np
import cv2
import matplotlib.pyplot as plt
 
img = cv2.imread("contoh.jpg", cv2.IMREAD_GRAYSCALE)
M, N = img.shape           # Tinggi, Lebar
 
# 1. Rotasi 30 derajat terhadap pusat citra
pusat = (N // 2, M // 2)
A_rot = cv2.getRotationMatrix2D(pusat, angle=30, scale=1.0)
img_rot = cv2.warpAffine(img, A_rot, (N, M))
 
# 2. Translasi: geser 50px ke kanan, 30px ke bawah
A_trans = np.float32([[1, 0, 50],
                       [0, 1, 30]])
img_trans = cv2.warpAffine(img, A_trans, (N, M))
 
# 3. Scaling: perbesar 1.5x
A_scale = np.float32([[1.5, 0,   0],
                       [0,   1.5, 0]])
img_scale = cv2.warpAffine(img, A_scale, (N, M))
 
# Tampilkan hasil
fig, axes = plt.subplots(1, 4, figsize=(14, 4))
for ax, im, title in zip(axes,
                          [img, img_rot, img_trans, img_scale],
                          ["Original", "Rotasi 30 deg",
                           "Translasi", "Scaling 1.5x"]):
    ax.imshow(im, cmap="gray")
    ax.set_title(title)
    ax.axis("off")
plt.tight_layout()
plt.show()
 
# Bonus: komposisi transformasi via perkalian matriks 3x3
theta  = np.radians(30)
A_rot3 = np.array([[ np.cos(theta), -np.sin(theta), 0],
                    [ np.sin(theta),  np.cos(theta), 0],
                    [ 0,              0,             1]])
A_tra3 = np.array([[1, 0, 50],
                    [0, 1, 30],
                    [0, 0,  1]])
 
# Rotasi LALU translasi -- cukup kalikan matriksnya
A_total = A_tra3 @ A_rot3
print("Matriks transformasi total (rotasi -> translasi):")
print(np.round(A_total, 4))
\end{lstlisting}
 
% ===============================================================================
\section*{Ringkasan}
\addcontentsline{toc}{section}{Ringkasan}
 
\begin{center}
\renewcommand{\arraystretch}{1.4}
\begin{tabular}{@{}lll@{}}
  \toprule
  \textbf{Topik} & \textbf{Konsep Kunci} & \textbf{Persamaan Utama} \\
  \midrule
  Persepsi Visual    & Adaptasi $B_a$, Mach Bands, Contrast  & -- \\
  Spektrum EM        & Gelombang \& foton                    & $\lambda\nu=c$,\ $E=h\nu$ \\
  Model Citra        & Iluminasi $\times$ Reflektansi         & $f=i\cdot r$ \\
  Sampling \& Kuant. & Resolusi spasial \& intensitas         & $b=M\times N\times k$ \\
  Ketertetanggaan    & $N_4$, $N_8$, $D_e$, $D_4$, $D_8$   & Pers.~\eqref{eq:eucl}--\eqref{eq:chess} \\
  Transformasi       & Linearitas, Affine, Koordinat Homogen & Pers.~\eqref{eq:linear},~\eqref{eq:affine} \\
  \bottomrule
\end{tabular}
\end{center}
 
\end{document}
 
